

\begin{table*}[t]
\centering
% \setlength{\tabcolsep}{6pt}
\caption{
\textbf{Evaluation on joint world-centric geometry and motion reconstruction.}
All metrics are reported without percentage symbols for readability.
% Lower values of $\text{Rel}^p$ and EPE indicate better accuracy, while higher $\delta^p$ and APD denote stronger performance. 
* denotes not zero-shot scene flow evaluation.
-S and -P denote the Sequence mode and Pair mode of ST4RTrack.
Since ST4RTrack always compares with the first frame for motion,
for a fair comparison, we run it on every pair of consecutive frames and then transform the results into the world coordinate system using VGGT poses.
Plus, we add Zero-MSF + GT pose as a reference.
}

% \vspace{-0.5em}
\label{tab:joint_geo_motion}
\begin{adjustbox}{width=\textwidth}
\begin{tabular}{@{}l cc cc cc cc cc c@{}}
\toprule
% \multirow{2}{*}{\textbf{Task}} & 
\multirow{1}{*}{\textbf{Method}} &
\multicolumn{2}{c}{\textbf{Kubric}~\cite{greff2021kubric}} &
\multicolumn{2}{c}{\textbf{Spring}~\cite{Mehl2023_Spring}} &
\multicolumn{2}{c}{\textbf{VKITTI2}~\cite{cabon2020virtual}} &
\multicolumn{2}{c}{\textbf{Dynamic Replica}~\cite{karaev2023dynamicstereo}} &
\multicolumn{2}{c}{\textbf{Point Odyssey}~\cite{zheng2023pointodyssey}} &
\textbf{Avg.} \\
\cmidrule(lr){2-3} \cmidrule(lr){4-5} \cmidrule(lr){6-7} \cmidrule(lr){8-9} \cmidrule(lr){10-11}
\ \underline{\emph{Geometry}} &
$\text{Rel}^{p}\!\downarrow$ & $\delta^{p}\!\uparrow$ &
$\text{Rel}^{p}\!\downarrow$ & $\delta^{p}\!\uparrow$ &
$\text{Rel}^{p}\!\downarrow$ & $\delta^{p}\!\uparrow$ &
$\text{Rel}^{p}\!\downarrow$ & $\delta^{p}\!\uparrow$ &
$\text{Rel}^{p}\!\downarrow$ & $\delta^{p}\!\uparrow$ &
\textbf{Rank} $\!\downarrow$ \\
\cmidrule(lr){2-12}
% \multirow{7}{*}{\centering\rotatebox[origin=c]{90}{Geometry}} 
% \ \underline{\emph{Geometry}} & & & & & & & & & & & \\
POMATO~\cite{zhang2025pomato} + VGGT & 25.56 & 77.85 & 98.13 & 61.71 & 32.16 & 54.89 & 7.26 & 94.64 & 19.88 & 80.08 & 5.0 \\
ST4RTrack-S~\cite{st4rtrack2025} + VGGT & 6.61 & 95.59 & 123.27 & 42.26 & 67.68 & 21.87 & 5.65 & 96.29 & 31.00 & 68.17 & 4.0 \\
ST4RTrack-P~\cite{st4rtrack2025} + VGGT & 17.81 & 80.76 & 157.05 & 38.00 & 84.77 & 14.46 & 4.87 & 97.13 & 29.13 & 71.66 & 3.4 \\
DELTA~\cite{ngo2024delta} + VGGT & 14.09 & 85.73 & 106.88 & 50.81 & 57.03 & 42.76 & 23.21 & 74.35 & 51.47 & 50.92 & 6.0 \\
Zero-MSF~\cite{liang2025zero} + VGGT & 8.79 & 94.73 & 142.44 & 40.66 & 15.76 & 80.04 & 7.11 & 97.03 & 22.55 & 78.27 & 4.6 \\
\textcolor{gray}{Zero-MSF~\cite{liang2025zero} + GT} 
& \textcolor{gray}{8.78} & \textcolor{gray}{94.73} 
& \textcolor{gray}{142.45} & \textcolor{gray}{40.69} 
& \textcolor{gray}{11.22} & \textcolor{gray}{89.92} 
& \textcolor{gray}{7.11} & \textcolor{gray}{97.01} 
& \textcolor{gray}{22.52} & \textcolor{gray}{78.27} 
& - \\
\method (ours) & \textbf{3.40} & \textbf{98.73} & \textbf{29.20} & \textbf{77.27} & \textbf{14.60} & \textbf{84.58} & \textbf{4.04} & \textbf{99.00} & \textbf{9.94} & \textbf{94.90} & \textbf{1.0} \\
\midrule
\midrule
\ \underline{\emph{Motion}} & 
EPE$\!\downarrow$ & APD$_{0.05}\!\uparrow$ &
EPE$\!\downarrow$ & APD$_{0.1}\!\uparrow$ &
EPE$\!\downarrow$ & APD$_{0.3}\!\uparrow$ &
EPE$\!\downarrow$ & APD$_{0.05}\!\uparrow$ &
EPE$\!\downarrow$ & APD$_{0.05}\!\uparrow$ &
\\
\cmidrule(lr){2-12}
% \multirow{7}{*}{\centering\rotatebox[origin=c]{90}{Motion}}
% \underline{\emph{Motion}} & & & & & & & & & & & \\
POMATO~\cite{zhang2025pomato} + VGGT & 79.58 & 5.23 & 180.13 & 20.16 & 368.78 & 15.66 & 14.56* & 36.34* & 31.94* & 30.32* & 5.0 \\
ST4RTrack-S~\cite{st4rtrack2025} + VGGT & 59.34* & 1.68* & 105.02 & 33.47 & 156.85 & 22.11 & 1.12* & 97.13* & 10.37* & 53.39* & 4.0 \\
ST4RTrack-P~\cite{st4rtrack2025} + VGGT & 217.20* & 5.65* & 441.84 & 9.07 & 874.94 & 13.16 & 0.99* & 97.55* & 58.62* & 51.19* & 4.8 \\
DELTA~\cite{ngo2024delta} + VGGT & 8.29* & 52.95* & 8.59 & 81.55 & 156.64 & 17.7 & 0.75 & 99.57 & 6.09 & 62.11 & 3.2 \\
Zero-MSF~\cite{liang2025zero} + VGGT & 8.59* & 50.13* & 7.78* & 85.59* & 112.99* & 21.69* & 0.59* & \textbf{99.80}* & 3.58* & 80.12* & 2.4 \\
\textcolor{gray}{Zero-MSF~\cite{liang2025zero} + GT} & \textcolor{gray}{5.74*} & \textcolor{gray}{72.37*} & \textcolor{gray}{5.50*} & \textcolor{gray}{87.59*} & \textcolor{gray}{73.81*} & \textcolor{gray}{25.15*} & \textcolor{gray}{0.40*} & \textcolor{gray}{99.80*} & \textcolor{gray}{2.30*} & \textcolor{gray}{91.04*} & - \\
\method (ours) & \textbf{4.60}* & \textbf{68.01}* & \textbf{5.61}* & \textbf{90.17}* & \textbf{71.75}* & \textbf{25.90}* & \textbf{0.51} & 99.72 & \textbf{3.49} & \textbf{80.66} & \textbf{1.0} \\
\bottomrule
\end{tabular}
\end{adjustbox}
% \vspace{-1em}
\end{table*}



\begin{figure*}
    \centering
    \includegraphics[width=\linewidth]{figs/indomain_result_figure.pdf}
    \caption{\textbf{Qualitative comparison with Zero-MSF~\cite{liang2025zero}.} Zoom in for the details. Compared to Zero-MSF, we have a more reasonable scene structure and better geometric details. More importantly, our predicted 3D scene flow has a more accurate direction of motion.}
    \label{fig:indomain_result}
    \vspace{-1em}
\end{figure*}

\section{Experiments}
\label{sec:experiments}
\subsection{Evaluation Setting}

\paragraph{Datasets}
For \emph{geometry} evaluation, we perform zero-shot testing on three unseen dynamic scene datasets:
DDAD~\cite{ddad2020},
Monkaa~\cite{mayer2016large},
and Sintel~\cite{sintel}. 
These datasets cover both real-world and synthetic scenes,
including indoor and outdoor environments. 
% 
For \emph{motion} evaluation,
due to the limited availability of datasets with dense scene flow annotations,
we use a combination of three in-domain datasets
(Kubric~\cite{greff2021kubric},
Spring~\cite{Mehl2023_Spring},
and VKITTI2~\cite{cabon2020virtual})
and two out-of-domain datasets
(Dynamic Replica~\cite{karaev2023dynamicstereo}
and Point Odyssey~\cite{zheng2023pointodyssey}). 
Since Dynamic Replica and Point Odyssey only provide sparse scene flow annotations,
we compute metrics only on the annotated points.

\paragraph{Metrics}
Unlike previous methods,
we evaluate geometry and motion in the \textit{world coordinate system}.
% to directly assess the real reconstruction ability of the model. 
The predicted world-space point cloud is aligned with the ground truth by optimizing per-sequence scale and shift parameters. 
We report the \textit{relative point error} ($\text{Rel}^p$) and the \textit{percentage of inlier} ($\delta^p$, threshold 0.25) as evaluation metrics. 
% 
The predicted scene flow is aligned according to the point map scale.
We compute the \textit{End Point Error} (EPE) and the \textit{Average Percent of Points within Delta} (APD), 
where the subscript of APD denotes the inlier threshold in the metric scale. 
We provide details in the supplementary material.

\subsection{Comparison with the State-of-the-art Methods}

\paragraph{Assessing joint Geometry and Motion Reconstruction}
We compare \method with recent representative methods in joint geometry and motion estimation in~\cref{tab:joint_geo_motion}.
Here,
we assess these metrics in the world coordinate space.
Our model directly outputs a sequence of predictions in world coordinates.
By contrast, most existing methods follow a pairwise design based on DUSt3R~\cite{wang2024dust3r},
which need post-optimization or camera poses to align with ground truth.
To ensure fairness, we use the camera pose predicted by VGGT~\cite{wang2025vggt} to transform their predictions into the world coordinate system.
% 
Through both quantitative and qualitative comparison in~\cref{fig:indomain_result,tab:joint_geo_motion},
we observe that these pairwise approaches usually exhibit degraded performance when extended to video sequences, while our method outperforms state-of-the-art methods by 38.64\% in geometry and 25.0\% in motion on average.
Note that,
unlike Zero-MSF~\cite{liang2025zero},
we do not train our model with motion annotations from Dynamic Replica~\cite{karaev2023dynamicstereo} and Point Odyssey~\cite{zheng2023pointodyssey}, yet still achieve better performance except on one comparable metric.
\Cref{fig:effectiveness} shows that our method estimates temporally consistent scene flow in a world coordinate system, 
remaining robust to camera motion and capable of describing 4D scene dynamics more accurately and efficiently.


\begin{table}[t]
\centering
\caption{
\textbf{Evaluation on world-centric geometric reconstruction.}
% Lower $\text{Rel}^p$ and higher $\delta^p$ indicate better performance.
$^\dagger$denotes using post-optimization. 
Note that,
our results are reported without any post-optimization.
}
\vspace{-0.5em}
\label{tab:geo_only}
\begin{adjustbox}{width=\linewidth}
\begin{tabular}{@{}l cc cc cc c@{}}
\toprule
\multirow{2}{*}{\textbf{Method}} &
\multicolumn{2}{c}{\textbf{Monkaa}~\cite{mayer2016large}} &
\multicolumn{2}{c}{\textbf{Sintel}~\cite{sintel}} &
\multicolumn{2}{c}{\textbf{DDAD}~\cite{ddad2020}} &
\textbf{Avg.}  \\
\cmidrule(lr){2-3} \cmidrule(lr){4-5} \cmidrule(lr){6-7}
&
$\text{Rel}^{p}\!\downarrow$ & $\delta^{p}\!\uparrow$ &
$\text{Rel}^{p}\!\downarrow$ & $\delta^{p}\!\uparrow$ &
$\text{Rel}^{p}\!\downarrow$ & $\delta^{p}\!\uparrow$ &
\textbf{Rank} $\!\downarrow$
\\
\midrule
\ \underline{Camera-centric} & & & & & & & \\ 
DepthPro~\cite{bochkovskiy2025depth} & 36.96 & 63.65 & 43.30 & 42.30 & 35.38 & 45.49 & 7.00 \\
MoGe~\cite{wang2024moge} & 35.21 & 65.95 & 35.28 & 60.97 & 18.63 & 77.07 & 4.33 \\
GeoCrafter~\cite{xu2025geometrycrafter} & 33.44 & 65.79 & 30.61 & 68.07 & 19.17 & 76.39 & 3.33 \\
\midrule
\ \underline{World-centric} & & & & & & & \\
MonST3R~\cite{zhang24monst3r}$^\dagger$ & 41.41 & 31.46 & 37.65 & 51.41 & 31.56 & 55.30 & 6.33 \\
VGGT~\cite{wang2025vggt} & 34.54 & 56.65 & \textbf{26.83} & \textbf{67.91} & 15.98 & 84.06 & \textbf{2.33} \\
Geo4D~\cite{jiang2025geo4d}$^\dagger$ & 28.04 & 69.52 & 34.61 & 59.54 & \textbf{14.58} & \textbf{83.68} & \textbf{2.33} \\
St4RTrack~\cite{st4rtrack2025} & 47.04 & 45.46 & 40.59 & 51.54 & 39.59 & 32.47 & 7.67 \\
\textbf{Ours} & \textbf{25.88} & \textbf{74.01} & 32.46 & 63.14 & 21.27 & 72.82 & 2.67 \\
\bottomrule
\end{tabular}
\end{adjustbox}
\vspace{-1em}
\end{table}


\paragraph{Assessing Geometry Reconstruction}
We further compare our method with several representative approaches in geometry reconstruction in~\cref{tab:geo_only}.
We group them by their geometric representation:
(1) Camera-centric methods predict depth or point maps in the camera coordinate system.
We transform their outputs to world coordinates using VGGT poses for fair comparison.
(2) World-centric methods predict geometry directly in the reference frame coordinate system, and we align their results to ground truth using an affine transformation.
% 
Our method achieves \emph{state-of-the-art} performance on Monkaa, 
demonstrating the advantages of our proposed architecture. 
For Sintel and DDAD, our performance is inferior to VGGT~\cite{wang2025vggt}, 
which we attribute to our single-modal design (w/o camera rays and depth maps) and the limited scale of our outdoor training data. 
Note that we do not perform post-optimization, as in Geo4D~\cite{jiang2025geo4d}.
The visual comparisons (see in supp.) further validate that our approach produces more coherent and consistent reconstructions in dynamic environments.

\begin{figure}
    \vspace{-0.5em}
    \centering
    \includegraphics[width=\linewidth]{figs/effectiveness_figure.pdf}
    \caption{\textbf{Qualitative comparison with ST4RTrack~\cite{st4rtrack2025}.}
    In the first case, the pixel trajectory shows that we yield cleaner scene flow, while ST4RTrack suffers from noisy drift.
    In the second case, the deformed point map (with darker color) shows that our method predicts more temporally consistent geometry and motion.
    }
    \label{fig:effectiveness}
    \vspace{-0.5em}
\end{figure}


\subsection{Ablation Study}
We conduct thorough ablations to analyze \method.
Results are reported in~\cref{tab:ablation_geometry,tab:ablation_motion},
aiming to investigate the following three key questions:

\begin{table}[t]
\centering
\caption{
\textbf{Ablation study on geometry VAE.}
Here,
we report geometry reconstruction results on both VAE and U-Net,
as a better VAE may not always lead to better final results.
The models are trained on a subset only for geometry reconstruction.
}
\label{tab:ablation_geometry}
\vspace{-0.5em}
\begin{adjustbox}{width=\linewidth}
\begin{tabular}{@{}l cc cc cc@{}}
\toprule
\multirow{2}{*}{\textbf{Model}} & 
\multirow{2}{*}{\textbf{Training Type}} & 
\multirow{2}{*}{\textbf{Rescale}} & 
\multicolumn{2}{c}{\textbf{Sintel}\cite{sintel}} &
\multicolumn{2}{c}{\textbf{Monkaa}~\cite{mayer2016large}} \\
\cmidrule(lr){4-5} \cmidrule(lr){6-7}
& & &
$\text{Rel}^{p}\!\downarrow$ & $\delta^{p}\!\uparrow$ &
$\text{Rel}^{p}\!\downarrow$ & $\delta^{p}\!\uparrow$ \\
\midrule
VAE - 1 & Original  & Max &
39.96 & 62.63 &
23.78 & 67.33 \\

VAE - 2 & From scratch & Max &
18.80 & 79.68 &
11.48 & 90.55 \\

VAE - 3  & Finetune decoder & Max &
20.76 & 77.77 &
14.44 & 85.91 \\

VAE - 4 & Finetune all & Mean &
\textbf{5.68} & \textbf{98.77} &
\textbf{5.03} & \textbf{99.13} \\

\midrule
Unet - I & VAE - 3 + Unet & Max &  40.47 & 49.42 & 33.66 & 56.42  \\
Unet - II & VAE - 4 + Unet & Mean &  \textbf{35.17} & \textbf{58.08} & \textbf{27.36} & \textbf{66.21} \\

\bottomrule
\end{tabular}
\end{adjustbox}
\vspace{-0.5em}
\end{table}



\noindent
\textbf{Is it necessary to strictly align the input distribution with that of Video Diffusion?}
The answer is \textbf{no}.
Most of the existing geometric diffusion models~\cite{ke2024repurposing,zhang2024world,jiang2025geo4d}
strictly normalize the 3D attributes (such as depth and point maps) to $[-1, 1]$,
in order to inherit the diffusion prior.
However,
we find that such a \textit{max-rescale} normalization results in suboptimal reconstruction accuracy for the pretrained VAEs (VAE-1\&2 of~\cref{tab:ablation_geometry}). 
Although Geo4D~\cite{jiang2025geo4d} alleviates this issue by freezing the VAE encoder and fine-tuning only the decoder, 
our ablation shows that this strategy still yields inferior results (VAE-3 in~\cref{tab:ablation_geometry}).
In contrast,
our proposed \emph{mean rescale} strategy---combined with fine-tuning all VAE components---achieves the best performance without strict adherence to the original VAE distribution (VAE-4 of~\cref{tab:ablation_geometry}). 
We further validate this finding in the Diffusion Unet stage (Unet-I vs. Unet-II),
resulting in an average of 16.6\% gain in geometry.
This finding suggests that \method maintains a better generalization ability,
even do \emph{not} align the distribution with that of video diffusion.


\begin{table}[t]
\centering
% \renewcommand{\arraystretch}{1.1}
\caption{\textbf{Ablation study on motion VAE}.
Comparison of different designs across three dynamic scene flow datasets.
Again, we report results on both VAE and U-Net.
}
\label{tab:ablation_motion}
\vspace{-0.5em}
\begin{adjustbox}{width=\linewidth}
\begin{tabular}{@{}l c cc cc@{}}
\toprule
\multirow{2}{*}{\textbf{Model}} & 
\multirow{2}{*}{\textbf{Fusion Type}} & 
\multicolumn{2}{c}{\textbf{Spring}~\cite{Mehl2023_Spring}} &
\multicolumn{2}{c}{\textbf{Point Odyssey}~\cite{zheng2023pointodyssey}} \\
\cmidrule(lr){3-4} \cmidrule(lr){5-6}
 &  & EPE & APD$_{0.03}\uparrow$ & EPE & APD$_{0.03}\uparrow$ \\
\midrule
VAE - 5 & Original &  6.43 & 50.49 & 1.55 & 91.68 \\
VAE - 6 & Offset &  1.83 & 83.51 & 2.00 & 86.66 \\
VAE - 7 & Separate &  \textbf{0.66} & \textbf{96.75} & \textbf{0.77} & \textbf{94.74} \\
VAE - 8 & Unify &  0.88 & 94.78 & \textbf{0.77} & 94.19 \\
\midrule
Unet -III & Separate & 6.37 & 65.94 & 3.65 & 63.19 \\
Unet - IV & Unify & \textbf{5.16} & \textbf{72.81} & \textbf{3.49} & \textbf{66.36} \\
\bottomrule
\end{tabular}
\end{adjustbox}
\vspace{-0.5em}
\end{table}


\noindent
\textbf{Which strategy for fusing geometric and motion latent information is most effective?}
We explore different strategies,
including \emph{Offset}, \emph{Separate}, and \emph{Unify} as discussed in~\cref{sec:method_vae},
for jointly encoding geometry and motion information in~\cref{tab:ablation_motion}.
Experimental results indicate that although separate VAEs achieve optimal performance in VAE reconstruction, a unified VAE ultimately performs better in Unet prediction.
This phenomenon highlights the importance of tightly coupling geometry and motion representations for coherent 4D modeling.

\noindent
\textbf{What prior knowledge does video generator provide?}
When using the original pretrained VAE to encode geometry and motion,
we observe that the model already exhibits reasonable reconstruction ability in indoor scenes, as seen in the first rows of~\cref{tab:ablation_geometry,tab:ablation_motion}. 
However, due to the significant scale discrepancy between point maps, scene flow, and image-space distributions,
the original VAE struggles to handle large-scale variations, especially in outdoor scenes, shown in~\cref{fig:rescale}. 
We argue that appropriate normalization and fine-tuning strategies are essential to fully leverage the priors embedded in the diffusion model. 
As shown in VAE-2 of Table~\ref{tab:ablation_geometry}, 
training from scratch leads to suboptimal results, confirming that the pretrained Video Diffusion model indeed contains rich priors beneficial for dense 4D reconstruction. 
By explicitly modeling these priors, we effectively unlock the 4D representation capability of the video diffusion model.




